% c07 = p024–p028（第 16–17 节；p024 顶部为第 15 节正文的续篇）
% [Tutor Note: p024 顶部承接 p023 末尾「群の準同型 φ: ℂ^× → K を (15.11) で定義する。」一句，
%  故本 chunk 以「来定义。」起句；第 15 节的续文到 (15.18) 与转置镜像对称的说明为止，
%  其后才是第 16 节标题。]

来定义。但 $\ell$ 是对所有 $i=1,\dots,n$ 使 $\ell\varphi_i$ 为整数的最小非负整数。此时 $\varphi$ 是单射，将其余核记作 $\overline{G}_{\max}$；

\begin{equation}
  1 \to \C^{\times} \xrightarrow{\ \varphi\ } K \to \overline{G}_{\max} \to 1. \tag{15.12}
\end{equation}

$W$ 通过

\begin{equation}
  \deg x_i = \ell\varphi_i,\quad i=1,\dots,n \tag{15.13}
\end{equation}

而具有 $\Z$-分次。由于 $f$ 在原点处有孤立临界点，加权超曲面

\begin{equation}
  Y := \bigl\{[x_1:\cdots:x_n]\in\PP(\ell\varphi_1,\dots,\ell\varphi_n) \bigm| f(x_1,\dots,x_n)=0\bigr\} \tag{15.14}
\end{equation}

是光滑的 Deligne–Mumford 栈。$Y$ 的典范丛成为平凡丛的充要条件是

\begin{equation}
  \deg x_1 + \cdots + \deg x_n = \ell \tag{15.15}
\end{equation}

成立。$\operatorname{Im}\varphi\cap G$ 由

\begin{equation}
  J := \bigl(\exp(2\pi\sqrt{-1}\varphi_1),\dots,\exp(2\pi\sqrt{-1}\varphi_n)\bigr) \tag{15.16}
\end{equation}

生成。假设 $J\in G$，令 $\overline{G} := G/\langle J\rangle$，并把自然满射 $K\to\overline{G}_{\max}$ 下 $\overline{G}$ 的原像记作 $M$。此外，把 $M$ 的特征标群记作 $H$。此时，Calabi–Yau 栈 $[Y/\overline{G}]$ 与 Landau–Ginzburg 轨形 $(f,G)$（或 $(f,M)$）\footnote{代数流形上的全纯函数 $f\colon V\to\C$ 与作用在 $V$ 上且使 $f$ 不变的群 $G$ 所成的组 $(f,G)$，这里称为 Landau–Ginzburg 轨形。}密切相关这一猜想称为 Calabi–Yau/Landau–Ginzburg 对应（Calabi–Yau/Landau–Ginzburg correspondence）。这至少是回溯到 [Gep87] 的古老想法，而最近的发展例如有 [Orl09, CR10, CR11, CIR14] 等。

可逆多项式 $f$ 的转置（transpose）由

\begin{equation}
  \check{f} := \sum_{i=1}^{n} x_1^{a_{1i}}\cdots x_n^{a_{ni}} \tag{15.17}
\end{equation}

定义。$\check{f}$ 的最大对角对称性群 $\check{G}_{\max}$ 由 $(A^T)^{-1}=(A^{-1})^T$ 的列向量 $\check{\rho}_i$ 生成。对 $G_{\max}$ 的子群 $G$，$G$ 的转置 $\check{G}$ 由

\begin{equation}
  \check{G} := \left\{\prod_{i=1}^{n}\check{\rho}_i^{r_i} \ \middle|\
    \text{对任意 } \prod_{i=1}^{n}\rho_i^{a_i}\in G \text{，均有 } (r_1\ \cdots\ r_n)\,A^{-1}
    \begin{pmatrix} a_1\\ \vdots\\ a_n\end{pmatrix}\in\Z \right\} \tag{15.18}
\end{equation}

定义（见 [Kra]）。对应 $G\leftrightarrow\check{G}$ 分别把 $G_{\max}$ 与 $\{e\}$、把 $\langle J\rangle$ 与 $G_{\max}\cap SL_n(\C)$ 互换。此时，Landau–Ginzburg 模型 $(f,G)$ 与 $(\check{f},\check{G})$ 互为镜像，这就是 Berglund–Hübsch [BH93] 的转置镜像对称（transposition mirror symmetry）。

\section{箭图及其表示}

同调镜像对称的证明方针大致分为两种。一种是用 Fourier 变换那样的几何方法定义函子，并证明它是满忠实且本质满射的；另一种则是，对每个范畴找出其生成元，使这些生成元的自同构代数作为微分分次代数或 A∞ 代数拟同构。这些方法类似于：要证明两个环同构，先作出同态，再证明它既满又单；也类似于在生成元之间建立对应，再证明关系式一致。

% [Tutor Note: 原文作「自己同型代数」，按上下文疑为「自己準同型代数」（自同态代数）之笔误，此处按原文译作「自同构代数」，请审校确认。]

给定预完备的增广三角范畴 $\cD$ 的紧生成元 $E$ 时，对于 $\cA := \operatorname{hom}_{\cD}(E,E)$，函子

\begin{equation}
  \operatorname{RHom}_{\cD}(E,-)\colon \cD \to D\cA \tag{16.1}
\end{equation}

及其伴随函子

\begin{equation}
  (-)\otimes_{\cA} E\colon D\cA \to \cD \tag{16.2}
\end{equation}

给出范畴等价\footnote{虽然也有人强调这与 Fourier–SYZ 变换（或族的 Floer 上同调）的类似性，但有一个显著差异：拉格朗日环面纤维化即使把所有纤维合起来也不生成导出范畴（在镜像一侧看，点层不用说导出范畴，就连 $K$ 群也不生成）。}。若 $\cA$ 是非循环的（acyclic，即对 $i\ne 0$ 有 $H^i(\cA)=0$），则 $D\cA$ 等价于 $A := H^0(\cA)$ 上模的导出范畴。

由顶点（vertex）的集合 $Q_0$、箭（arrow）的集合 $Q_1$，以及对每个箭赋予其起点（source）与终点（target）的映射 $s,t\colon Q_1\to Q_0$ 组成的组 $Q := (Q_0,Q_1,s,t)$ 称为箭图（quiver）。当 $Q_0$ 与 $Q_1$ 都是有限集时，称 $Q$ 是有限的（finite）。对正整数 $k$，长度（length）为 $k$ 的路（path）是指箭的序列 $(a_k,\dots,a_2,a_1)$，且对任意 $i=1,\dots,k-1$ 满足 $s(a_{i+1})=t(a_i)$。另外，对每个顶点 $v\in Q_0$，考虑以 $v$ 为起点、以 $v$ 为终点、长度为 $0$ 的元素 $e_v$，这也称为路。在由路张成的向量空间上赋予乘法

\begin{equation}
  (a_k,\dots,a_1)\cdot(b_l,\dots,b_1) =
  \begin{cases}
    (a_k,\dots,a_1,b_l,\dots,b_1) & s(a_1)=t(b_l),\\
    0 & \text{其他}
  \end{cases} \tag{16.3}
\end{equation}

所得的代数称为路代数（path algebra），记作 $\C Q$。附着于顶点的长度为 $0$ 的路是这个代数的幂等元。路代数 $\C Q$ 具有由路长给出的自然分次 $\C Q=\bigoplus_{n=0}^{\infty}(\C Q)_n$。箭图 $Q$ 与其路代数的双边理想 $I\subset\C Q$ 组成的组 $\Gamma := (Q,I)$ 称为带关系箭图（quiver with relations），$\C\Gamma := \C Q/I$ 称为它的路代数。以下假设箭图的关系式都含于由长度 $2$ 以上的元素构成的子模中；

\begin{equation}
  I \subset (\C Q)_{\ge 2} := \bigoplus_{n=2}^{\infty}(\C Q)_n. \tag{16.4}
\end{equation}

对带关系箭图 $\Gamma$，其路代数 $\C\Gamma$ 上的模称为 $\Gamma$ 的表示（representation）。给出带关系箭图的表示，等价于给出：对应各顶点 $v$ 的线性空间 $M_v$ 与对应各箭 $a$ 的线性映射 $\varphi_a\in\Hom(M_{s(a)},M_{t(a)})$ 所成的组 $\bigl((M_v)_{v\in Q_0},(\varphi_a)_{a\in Q_1}\bigr)$，且满足关系式。

把由对应顶点的幂等元生成的半单代数记作

\begin{equation}
  \C^{Q_0} := \bigoplus_{v\in Q_0}\C e_v \tag{16.5}
\end{equation}

则带关系箭图的路代数是半单代数 $\C^{Q_0}$ 上的代数。反之，给定半单代数 $\C^{Q_0}$ 上的代数 $A$ 时，存在线性空间的直和分解 $A=\bigoplus_{v,w\in Q_0}e_v A e_w$，于是从各直和项中适当选取 $A$ 的生成元并置为 $Q_1$，当 $a\in Q_1$ 是 $e_v A e_w$ 的元素时定义 $s(a):=w$、$t(a):=v$，便得到箭图 $Q := (Q_0,Q_1,s,t)$。由构造，存在从 $\C Q$ 到 $A$ 的自然满射，取其核为 $I$ 并令 $\Gamma := (Q,I)$，则 $\C\Gamma$ 与 $A$ 自然同构。

对非循环的生成元 $E=E_1\oplus\cdots\oplus E_n$，准备与直和项个数相同数目的顶点并置 $Q_0 := \{v_1,\dots,v_n\}$，则通过把 $e_i := e_{v_i}$（$i=1,\dots,n$）送到 $\id_{E_i}\subset A$，$E$ 的自同态代数 $A$ 成为半单代数 $\C^{Q_0}$ 上的代数，从而可表示为适当的带关系箭图的路代数。

设 $\Gamma$ 为带关系箭图，$A := \C\Gamma$ 为其路代数。对任意顶点 $v\in Q_0$，把幂等元 $e_v\in A$ 生成的 $A$ 的右子模记作 $P_v := e_v A$。由 $1=\sum_{v\in Q_0}e_v$ 得 $A\cong\bigoplus_{v\in Q_0}P_v$，因此 $P_v$ 是 $A$ 的直和项，从而是射影模。当带关系箭图来自非循环的生成元 $E=E_1\oplus\cdots\oplus E_n$ 时，由范畴等价 (16.1)，$E_i$ 被送到 $P_i := P_{v_i}$。另一方面，给定顶点 $v$ 时，对顶点 $w$ 把线性空间定义为

\begin{equation}
  M_w =
  \begin{cases}
    \C & w=v,\\
    0 & \text{其他}
  \end{cases} \tag{16.6}
\end{equation}

并对所有箭 $a\in Q_1$ 令 $\varphi_a=0$，则 $\bigl((M_w)_{w\in Q_0},(\varphi_a)_{a\in Q_1}\bigr)$ 是 $\Gamma$ 的表示；把它记作 $S_v$，称为附着于 $v$ 的单模。$A$-模 $S_v$ 的射影分解的首项为

\begin{equation}
  \cdots \to \bigoplus_{t(a)=v}P_{s(a)} \to P_v \to S_v \to 0 \tag{16.7}
\end{equation}

给出，其核来自 $A$ 的关系式。把这个射影分解用函子 $\Hom(-,S_w)$ 送出，便可看出：$\Ext^1(S_v,S_w)$ 的基对应于从 $w$ 到 $v$ 的箭，$\Ext^2(S_v,S_w)$ 的基对应于从 $w$ 到 $v$ 的路的关系式。

不具有非平凡直和分解的对象称为不可分解的（indecomposable），不具有平凡子对象的对象称为不可约的（irreducible），除恒等映射的常数倍以外没有自同态的对象称为单的（simple）。又，单且非循环的对象称为例外对象（exceptional object）。若例外对象的序列 $(E_i)_{i=1}^n$ 对任意 $i>j$ 与任意 $k\in\Z$ 满足 $\Ext^k(E_i,E_j)=0$，则称为例外序列（exceptional collection）。非循环的例外序列称为强例外序列（strong exceptional collection）。$\bigoplus_{i=1}^n E_i$ 为生成元的例外序列称为完备例外序列（full exceptional collection）。

起点与终点一致的路称为有向循环（oriented cycle）。又，不含任何有向循环的箭图称为有向箭图（directed quiver）。完备强例外序列的全同态代数（total morphism algebra）$A := \bigoplus_{i,j=0}^n \Hom(E_i,E_j)\cong\End\bigl(\bigoplus_{i=0}^n E_i\bigr)$ 同构于带关系有向箭图的路代数。

% [Tutor Note: 原文作 ⊕_{i,j=0}^n Hom(E_i,E_j) ≅ End(⊕_{i=0}^n E_i)，而例外序列 (E_i)_{i=1}^n 的指标自 1 起，疑「0」应为「1」，此处按原文保留。]

在有向箭图的表示中，单模序列 $(S_n,S_{n-1},\dots,S_1)$ 构成完备例外序列，并且

\begin{equation}
  \Ext^k(P_i,S_j) =
  \begin{cases}
    \C & i=j \text{ 且 } k=0,\\
    0 & \text{其他}
  \end{cases} \tag{16.8}
\end{equation}

在这个意义上，它与射影模组成的强完备例外序列 $(P_i)_{i=1}^n$ 互为对偶。完备例外序列 $(S_i)_{i=n}^1$ 绝不会成为强例外序列，但若 $A$ 是 Koszul 代数（Koszul algebra），则在导出范畴中平移次数所得的序列 $(S_i[-i])_{i=n}^1$ 是强例外序列，且其全同态代数为 $A$ 的 Koszul 对偶代数（Koszul dual algebra）。在这个意义上，$(S_i[i])_{i=n}^1$ 的作为微分分次代数或 A∞ 代数的全同态代数，给出 $A$ 的 Koszul 对偶代数的推广。

作为箭图的例子，考虑由

\begin{equation}
  v_1 \xleftarrow{a_1} v_2 \xleftarrow{a_2} \cdots \xleftarrow{a_{n-1}} v_n \tag{16.9}
\end{equation}

给出的箭图 $Q$。它有 $n$ 个顶点 $Q_0=\{v_1,\dots,v_n\}$ 与 $n-1$ 支箭 $Q_1=\{a_1,\dots,a_{n-1}\}$，且对任意 $i=1,\dots,n-1$ 有 $s(a_i)=v_{i+1}$ 且 $t(a_i)=v_i$。对 $i,j=1,\dots,n$ 定义路 $a_{ij}$ 为

\begin{equation}
  a_{ij} :=
  \begin{cases}
    a_i a_{i+1}\cdots a_{j-1} & i<j,\\
    e_{v_i} & i=j,\\
    0 & i>j
  \end{cases} \tag{16.10}
\end{equation}

则路代数中的乘法由

\begin{equation}
  a_{ij}a_{kl} =
  \begin{cases}
    a_{il} & j=k,\\
    0 & \text{其他}
  \end{cases} \tag{16.11}
\end{equation}

给出。把除 $(i,j)$ 分量外全为 $0$、该分量为 $1$ 的 $n\times n$ 矩阵记作 $E_{ij}$，则通过映射 $a_{ij}\mapsto E_{ij}$，$Q$ 的路代数 $A := \C Q$ 同构于 $n$ 次上三角矩阵代数。由射影模 $P_i := e_i A$ 组成的强完备例外序列 $(P_n,P_{n-1},\dots,P_1)$ 的对偶序列由单模 $S_i := \C e_i$ 组成，而它们之间射构成的线性空间为

\begin{equation}
  \Ext^k(S_i,S_j) \cong
  \begin{cases}
    \C\cdot\id_{S_i} & i=j \text{ 且 } k=0,\\
    \C & j=i+1 \text{ 且 } k=1,\\
    0 & \text{其他}
  \end{cases} \tag{16.12}
\end{equation}

分次代数 $A := \bigoplus_{k=0}^{\infty}\bigoplus_{i,j=1}^n \Ext^k(S_i,S_j)$ 的乘法结构，容易看出由作为分次线性空间的结构 (16.12) 唯一确定。进一步，微分分次代数 $\cA$ 是形式的这一事实也容易由 (16.12) 推出。一般地，分次代数 $A$ 的增广\footnote{对给定的分次代数 $A$，上同调代数为与 $A$ 同构的微分分次代数 $\cA$ 称为 $A$ 的增广（enhancement）。} 唯一的现象称为内在形式性（intrinsic formality）。

把 Fukaya–Seidel 范畴中的射空间 (13.12) 与矩阵因子化范畴中的 (14.2)，同 (16.12) 所给分次代数的内在形式性结合起来，便对 $f := x^{n+1}-ax$ 与 $\check{f} := x^{n+1}$ 得到增广三角范畴的等价

\begin{equation}
  D^b\cF(f) \cong D^b_{\mathrm{sing}}\bigl(\operatorname{gr}\C[x]/(\check{f})\bigr) \tag{16.13}
\end{equation}

这就是关于 $A_n$ 奇点的同调镜像对称。类似的结果被猜想对更一般的可逆多项式也成立，并已在若干情形得到证明；关于这一点，例如可参见 [FU]。

\section{环面 Fano 流形的同调镜像对称}

在关于 Calabi–Yau 流形以外的流形的同调镜像对称中，研究得最深入的是环面 Fano 流形的情形。环面 Fano 流形的镜像是 Landau–Ginzburg 模型；由于这一明显的不对称性，按在哪一侧考虑辛几何而有两个版本。

环面 Fano 流形 $X$ 由凸格多面体\footnote{有限个格点的凸包（convex hull）称为凸格多面体（convex lattice polytope）。}决定，其复结构没有模空间。另一方面，辛结构在把 $X$ 构造为辛向量空间在环面作用下的辛商时，会随动量映射的值（等位集）的选取而变化，因而具有 $b_2=\dim H^2(X,\R)$ 维的模空间。这个模空间在镜像的 Landau–Ginzburg 模型中对应于势 $W$ 的系数。

同调镜像对称的一个版本是

\begin{equation}
  D^b\operatorname{coh}X \cong D^b\cF(W) \tag{17.1}
\end{equation}

而另一个版本是

\begin{equation}
  D^{\pi}\cF(X) \cong D^b_{\mathrm{sing}}(W) \tag{17.2}
\end{equation}

在无模空间依赖性这个意义上，(17.1) 更容易；但对一般的辛结构，期望 (17.2) 中出现的两个范畴都是半单的（即成为向量空间复形范畴的直和），在这个意义上 (17.2) 更容易。不过，给出 $X$ 的 Gromov–Witten 不变量与势 $W$ 的振荡积分之间关系、并应作为推论导出的经典镜像对称，其主张属于 (17.2) 一方。对未必是 Fano 流形的一般紧辛环面流形，(17.2) 的证明已由 Abouzaid–Fukaya–Oh–Ohta–Ono 预告。此外，关于此情形下同调镜像对称与经典镜像对称的关系，在 [FOOO] 中有详细讨论。不过，即便在这种情况下，也尚未从增广三角范畴的等价导出经典镜像对称。

以下讨论 (17.1)。

设 $\bm{N}$ 为秩 $n$ 的自由 Abel 群，其对偶 Abel 群记作 $\bm{M} := \Hom(\bm{N},\Z)$。设 $\Delta\subset\bm{N}_{\R} := \bm{N}\otimes\R$ 为凸格多面体，并假设 $\Delta$ 以原点为内点。把 $\Delta$ 的格点集记作 $A := \Delta\cap\bm{N}$，并把以 $\Delta$ 为 Newton 多面体的洛朗多项式记作

\begin{equation}
  W = \sum_{n\in A}a_n x^n \tag{17.3}
\end{equation}

这里若系数 $(a_n)_{n\in A}\in\C^A$ 取得足够一般，$W$ 所定义的映射 $\bm{M}^{\times}_{\C} := \Spec\C[\bm{N}]\to\C$ 给出 Lefschetz 纤维化。(17.1) 的右边就是如此得到的 Lefschetz 纤维化的 Fukaya–Seidel 范畴。但定义域 $\bm{M}^{\times}_{\C}\cong(\C^{\times})^n$ 上要赋予以 $-\sqrt{-1}/2\sum_{i=1}^n \dd x_i\wedge \dd\bar{x}_i/|x_i|^2$ 为辛形式的自然辛结构。

把 $\Delta$ 的顶点集记作 $\{n_i\}_{i=1}^r$，并定义映射 $\phi\colon\Z^r\to\bm{N}$ 为 $\phi(e_i) := n_i$。其中把 $\Z^r$ 的标准基记作 $(e_i)_{i=1}^r$。进一步，选取并固定 $\Delta$ 的边界（即余维 $1$ 的面的并集）的一个三角剖分 $\Sigma$，使其顶点集与 $\Delta$ 的顶点集一致。若存在某 $\{i_1,\dots,i_k\}\subset\{1,\dots,r\}$，使得 $z_{i_1}=\cdots=z_{i_k}=0$，但由 $\{v_{i_1},\dots,v_{i_k}\}$ 张成的单纯形不属于 $\Sigma$，这样的点 $(z_1,\dots,z_r)\in\C^r$ 的集合称为 Stanley–Reisner 轨迹（Stanley–Reisner locus），记作 $\mathrm{SR}(\Sigma)$。

% [Tutor Note: 原文此处作 {v_{i_1},…,v_{i_k}}，但上文把 Δ 的顶点记为 n_i，疑为记号混用或笔误；此处按原文保留 v。]

$\phi$ 诱导的映射 $\phi_{\C^{\times}}\colon(\C^{\times})^r\to\bm{N}_{\C^{\times}}$ 的核 $K := \operatorname{Ker}\phi_{\C^{\times}}\subset(\C^{\times})^r$ 通过 $(\C^{\times})^r$ 对 $\C^r$ 的自然作用而作用于 $\C^r$；关于这个作用，把 $\C^r\setminus\mathrm{SR}(\Sigma)$ 的栈论意义下的商记作

\begin{equation}
  X_{\Sigma} := [(\C^r\setminus\mathrm{SR}(\Sigma))/K] \tag{17.4}
\end{equation}

并称为附着于 $\Sigma$ 的环面栈。此时，猜想 2.5 断言增广三角范畴的等价

\begin{equation}
  D^b\operatorname{coh}X_{\Sigma} \cong D^b\cF(W) \tag{17.5}
\end{equation}

的存在。$X_{\Sigma}$ 一般是环面弱 Fano 栈\footnote{典范丛为 nef 且丰富的栈称为弱 Fano 栈。}。$X_{\Sigma}$ 成为代数流形的充要条件是 $\Sigma$ 为幺模的（即 $\Sigma$ 的任意单纯形都可通过 $GL_n(\Z)\ltimes\Z^n$ 的作用移到标准单纯形），而 $X_{\Sigma}$ 成为 Fano 的充要条件是 $\Sigma$ 的所有面都是 $\Delta$ 的面。$X_{\Sigma}$ 的同构类依赖于单纯剖分 $\Sigma$ 的选取，但导出范畴 $D^b\operatorname{coh}X_{\Sigma}$ 作为增广三角范畴的拟同构类仅由 $\Delta$ 决定 [Kaw05]。

作为例子，考虑 $\bm{N}:=\Z$ 是秩 $1$ 的 Abel 群、$\Delta_{a,b} := [-a,b]$ 是长度为 $a+b$ 的闭区间的情形。此时

\begin{equation}
  \phi_{\C^{\times}}\colon(\C^{\times})^2\to\C^{\times},\quad
  (\alpha,\beta)\mapsto\alpha^{-a}\beta^{b} \tag{17.6}
\end{equation}

而其核同构于 $\C^{\times}\times\Z/c\Z$，其中 $c := \gcd(a,b)$。$\partial\Delta$ 的单纯剖分是唯一的，Stanley–Reisner 轨迹为 $\bm{0} := \{(0,0)\}$。$K$ 在 $\C^2\setminus\bm{0}$ 上的作用在 $(1,0)$ 与 $(0,1)$ 处分别有同构于 $\Z/a\Z$ 与 $\Z/b\Z$ 的稳定化子。对应的环面栈记作

\begin{equation}
  \mathbb{X}_{a,b} := [(\C^2\setminus\bm{0})/K] \tag{17.7}
\end{equation}

若 $a$ 与 $b$ 互素，则 $\mathbb{X}_{a,b}$ 同构于 $\C^{\times}$ 通过作用 $(x,y)\mapsto(\alpha^a x,\alpha^b y)$ 在 $\C^2\setminus\bm{0}$ 上的商 $\PP(a,b) := [(\C^2\setminus\bm{0})/\C^{\times}]$。作为镜像，取

\begin{equation}
  W(x) = x^a - \frac{1}{x^b} \tag{17.8}
\end{equation}

则临界点由 $x^{a+b}=b/a$ 定义，其中一个是正实数。取 $\infty$ 为起点，选取连结临界值与 $\infty$ 的实轴上的半直线 $\gamma_1$ 作为消没路径，则对应的 Lefschetz 穿透 $L_1$ 成为 $x$ 平面的正实轴。令 $n=a+b$、$\zeta_n=\exp(2\pi\sqrt{-1}/n)$，对 $i\in\{1,\dots,n-1\}$，把用映射 $x\mapsto\zeta_n^i x$ 移动 $L_1$ 所得的记为 $\zeta_n^i L_1$，则它是沿 $\zeta_n^{ai}\gamma_1$ 的 Lefschetz 穿透。在紧集之外把 $\zeta_n^{ai}\gamma_1$ 形变为与正实轴平行，则 $\zeta_n^i L_1$ 在无穷远附近被旋转 $-\frac{2ai\pi}{n}$，在 $0$ 附近被旋转 $\frac{2bi\pi}{n}$。把这样得到的 Lefschetz 穿透记作 $L_{i+1}$，则 $(L_i)_{i=1}^n$ 成为 Lefschetz 穿透的特殊基。关于这个特殊基的 Fukaya–Seidel 范畴容易计算，若适当选取由 $\mathbb{X}$ 的线丛组成的完备强例外序列 $(\cL_i)_{i=1}^n$，则容易看出

% [Tutor Note: 原文此处作「Lefscheta 指貫」，疑为「Lefschetz 指貫」之笔误，正文按「Lefschetz 穿透」译出。]

\begin{equation}
  \Hom_{\mathfrak{Fuk}\,W}(L_i,L_j) \cong \Hom_{\mathbb{X}}(\cL_i,\cL_j) \tag{17.9}
\end{equation}

成立\footnote{若 $a$ 与 $b$ 互素，则可取 $\cL_i=\cO(i)$。若 $a$ 与 $b$ 不互素，则可描述为 $a$ 与 $b$ 互素情形的覆叠。}。

作为例子，考虑 $(a,b)=(1,1),(2,2),(3,5)$ 的情形。Lefschetz 穿透的特殊基分别如图 17.1、图 17.2 及图 17.3 所示，对应的箭图分别由 (17.10)、(17.11) 及 (17.12) 给出。注意在 $(a,b)=(1,1)$ 的情形 $\mathbb{X}=\PP^1$，而在 $(a,b)=(2,2)$ 的情形 $\mathbb{X}$ 是它的覆叠。

\begin{equation}
  1 \rightrightarrows 2 \tag{17.10}
\end{equation}

% [Tutor Note: (17.10) 之后的箭图 (17.11)、(17.12) 与图 17.1–17.3 位于 p029（属 c08），本 chunk 至 p028 末的 (17.10) 为止。]

% tags: (15.12)–(15.18) 7 条 + (16.1)–(16.13) 13 条 + (17.1)–(17.10) 10 条 = 30 条
% CHUNK c07 DONE pages p024–p028
